% !TEX TS-program = pdflatex
% !TeX program = pdflatex
% !TEX encoding = UTF-8
% !TEX spellcheck = en_US

\documentclass[xcolor=table]{beamer}

\usepackage{../extra/beamer/karimnlp}


%\usepackage{fullpage}
%\usepackage[left=2.8cm,right=2.2cm,top=2 cm,bottom=2 cm]{geometry}
\setbeamersize{text margin left=10pt,text margin right=10pt}



%\usepackage[T1]{fontenc}
\usepackage{textcomp}
\usepackage[utf8]{inputenc}
%\usepackage[french,english]{babel}
\usepackage{arabtex}
\usepackage{txfonts}
\usepackage{tipa}
\usepackage[]{graphicx}
\usepackage{multirow}
\usepackage{hyperref}
\usepackage{colortbl}
\usepackage{tabularray}
\usepackage{listingsutf8}
\usepackage{wrapfig}
\usepackage{multicol}
\usepackage[export]{adjustbox} %for images in table, also for frame
\usepackage[many]{tcolorbox}
\usepackage{wasysym}
\usepackage[lined]{algorithm2e}
\usepackage{alltt} %verbatim with commands
\usepackage{longtable}
\usepackage{tabu}

\usepackage{CJKutf8}

\usepackage{amsmath,amssymb, amsfonts} 

%\usepackage{newtxtext,newtxmath}


%\usepackage{acolor}


\definecolor{lightblue}{HTML}{D0D2FF}
\definecolor{lightyellow}{HTML}{FFFFAA}
\definecolor{darkblue}{HTML}{0000BB}
\definecolor{olivegreen}{HTML}{006600}
\definecolor{darkgreen}{HTML}{008B45} %009B55
\definecolor{violet}{HTML}{6600CC}
\definecolor{deeppink}{HTML}{FF1493}
\definecolor{orangey}{HTML}{FFBB00}


\makeatletter
\def\beamer@calltheme#1#2#3{%
	\def\beamer@themelist{#2}
	\@for\beamer@themename:=\beamer@themelist\do
	{\IfFileExists{\beamer@themelocation/#3\beamer@themename.sty}{%
	\usepackage[{#1}]{\beamer@themelocation/#3\beamer@themename}}{\usepackage[{#1}]{#3\beamer@themename}}
		
	}%
}

\def\usefolder#1{
	\def\beamer@themelocation{#1}
}
\def\beamer@themelocation{}
\makeatother

\usefolder{../../extra/beamer-themes}

\usetheme{Karim} % Antibes Boadilla Warsaw Karim

\beamertemplatenavigationsymbolsempty

\let\oldcite\cite
\renewcommand{\cite}[1]{{\bfseries\color{blue}\oldcite{#1}}}

\tcbuselibrary{listings}


%\renewcommand{\baselinestretch}{1.5}

\def\supit#1{\raisebox{0.8ex}{\small\it #1}\hspace{0.05em}}

\AtBeginSection{%
	\begin{frame}
		\sectionpage
	\end{frame}
}

\newcommand{\rottext}[2]{%
	\rotatebox{90}{%
	\begin{minipage}{#1}%
		\raggedleft#2%
	\end{minipage}%
	}%
}


\institute{ %
Laboratoire de la Communication dans les Systèmes Informatiques (LCSI)
\par
École  nationale Supérieure d'Informatique (ESI), Algiers, Algeria
}
\author[Abdelkrime Aries (2023/2024)] %
{Abdelkrime Aries}
%\titlegraphic{\includegraphics[height=1cm]{../img/esi-logo.png}%\hspace*{4.75cm}~


\date{Academic year: 2023/2024} %\today

\titlegraphic{%
	\includegraphics[height=1cm]{../../img/misc/esi-logo.png}%
	\hspace{2cm}%
	\includegraphics[height=1cm]{../../img/misc/lcsi-logo.png}%
	\hspace{2cm}%
	\includegraphics[height=1.5cm]{../../img/misc/esi.nlp-logo.png}%
}

%\setbeamertemplate{headline}{}

\newcommand{\kurl}[1]{{\scriptsize\bfseries\color{orangey}\url{#1}}}

\newcommand{\keyword}[1]{\textcolor{red}{\bfseries\itshape #1}}
\newcommand{\expword}[1]{\textcolor{olivegreen}{#1}}
\newcommand{\optword}[1]{\textcolor{violet}{\bfseries #1}}

\makeatletter
\newcommand\mysphere{%
	\parbox[t]{10pt}{\raisebox{0.2pt}{\beamer@usesphere{item projected}{bigsphere}}}}
\makeatother

%\let\oldtabular\tabular
%\let\endoldtabular\endtabular
%\renewenvironment{tabular}{\rowcolors{2}{white}{lightblue}\oldtabular\rowcolor{blue}}{\endoldtabular}


%\NoAutoSpacing %french autospacing after ":"

\def\graphpath{}

\newcommand{\changegraphpath}[1]{\def\graphpath{#1}}


\newcommand{\vgraphpage}[2][.84\textheight]{%
%	\begin{center}%
		\includegraphics[height=#1]{\graphpath #2}%
%	\end{center}%
}

\newcommand{\hgraphpage}[2][\textwidth]{%
%	\begin{center}%
		\includegraphics[width=#1]{\graphpath #2}%
%	\end{center}%
}

\newcommand{\graphpage}[2][]{%
	\includegraphics[#1]{\graphpath #2}%
}

\bibliographystyle{apalike}

\newcommand{\insertbibliography}[2]{
	\appendix
	\section*{Bibliography}
	\nocite{#2}
%	\makeatletter % to change template
%	\setbeamertemplate{headline}[default] % not mandatory, but I though it was better to set it blank
%	\def\beamer@entrycode{\vspace*{-\headheight}} % here is the part we are interested in :)
%	\makeatother
	\begin{multicols*}{2}[\frametitle{\insertsection} \usebeamertemplate*{frametitle}]%\usebeamertemplate*{frametitle}\frametitle{Références}
		\tiny
		\bibliography{#1}
	\end{multicols*}
}

\definecolor{my-grey}{RGB}{233, 233, 233}

\newcommand{\insertlicence}{
	\begin{frame}[plain]
	\frametitle{License: CC-BY 4.0}

	\begin{tcolorbox}[colback=cyan,
		colframe=cyan,  
		arc=0pt,outer arc=0pt,
		valign=top, 
		halign=center,
		width=\textwidth]
		
		\includegraphics[width=.5cm]{../../img/licence/cc_icon_white_x2.png}
		\includegraphics[width=.5cm]{../../img/licence/attribution_icon_white_x2.png}
		
		\color{white}
		\bfseries Attribution 4.0 International (CC BY 4.0) \\
		\tiny \url{https://creativecommons.org/licenses/by/4.0/deed.en}
		
	\end{tcolorbox}\vspace{-.5cm}
	\begin{tcolorbox}[colback=my-grey,
		colframe=my-grey,  
		center, arc=0pt,outer arc=0pt,
		valign=top, 
		halign=left,
		width=\textwidth]
		
		\tiny
		
		\begin{center}
			\bfseries\Large
			You are free to:
		\end{center}
		
		\begin{minipage}{0.83\textwidth}
			\begin{itemize}
				\item[] \textbf{Share} — copy and redistribute the material in any medium or format
				\item[] \textbf{Adapt} — remix, transform, and build upon the material
				for any purpose, even commercially.
			\end{itemize}
		
			\vspace{12pt}
		
			\textit{The licensor cannot revoke these freedoms as long as you follow the license terms.}
		\end{minipage}
		\begin{minipage}{0.15\textwidth}
			\includegraphics[width=\textwidth]{../../img/licence/FreeCulturalWorks_seal_x2.jpg}
		\end{minipage}
	
		
		\begin{center}
			\bfseries\Large
			Under the following terms:
		\end{center}
		
		\begin{itemize}
			\item[] \textbf{Attribution} — You must give appropriate credit, provide a link to the license, and indicate if changes were made. You may do so in any reasonable manner, but not in any way that suggests the licensor endorses you or your use.
			\item[] \textbf{No additional restrictions} — You may not apply legal terms or technological measures that legally restrict others from doing anything the license permits.
		\end{itemize}
		
	\end{tcolorbox}
	
%	\begin{center}
%		\bfseries Attribution 4.0 International (CC BY 4.0)
%		\url{https://creativecommons.org/licenses/by/4.0/deed.fr}
%	\end{center}

%	\tiny
%
%	Vous êtes autorisé à : 
%	\begin{itemize}
%		\item \textbf{Partager} — copier, distribuer et communiquer le matériel par tous moyens et sous tous formats
%		\item \textbf{Adapter} — remixer, transformer et créer à partir du matériel
%	\end{itemize}
%	
%	Selon les conditions suivantes : 
%	\begin{itemize}
%		\item \textbf{Attribution} — Vous devez créditer l'Œuvre, intégrer un lien vers la licence et indiquer si des modifications ont été effectuées à l'Oeuvre. Vous devez indiquer ces informations par tous les moyens raisonnables, sans toutefois suggérer que l'Offrant vous soutient ou soutient la façon dont vous avez utilisé son Oeuvre.
%		\item \textbf{Pas d'Utilisation Commerciale} — Vous n'êtes pas autorisé à faire un usage commercial de cette Oeuvre, tout ou partie du matériel la composant. 
%		\item \textbf{Pas de restrictions complémentaires} — Vous n'êtes pas autorisé à appliquer des conditions légales ou des mesures techniques qui restreindraient légalement autrui à utiliser l'Oeuvre dans les conditions décrites par la licence.
%	\end{itemize}

	\end{frame}
}

\settowidth{\leftmargini}{\usebeamertemplate{itemize item}}
\addtolength{\leftmargini}{\labelsep}

\AtBeginDocument{
	\newcolumntype{L}[2]{>{\vbox to #2\bgroup\vfill\flushleft}p{#1}<{\egroup}} 
	
	\begin{frame}[plain]
		\maketitle
	\end{frame}

	\insertlicence
}


% needs etoolbox; to break links after -
\appto\UrlBreaks{\do\-}


\makeatletter
\newcommand{\xRightarrow}[2][]{\ext@arrow 0359\Rightarrowfill@{#1}{#2}}
\makeatother


\usefonttheme{structurebold}
%\usefonttheme{professionalfonts}


% figure caption
\setlength\abovecaptionskip{2pt}

% equations
\AtBeginDocument{
\abovedisplayskip=0pt
\abovedisplayshortskip=0pt
\belowdisplayskip=0pt
\belowdisplayshortskip=0pt
}


\subtitle[03- Basic text processing]{Chapter 03\\Basic text processing} 

\changegraphpath{../img/basic/}

\begin{document}
	
\begin{frame}
	\frametitle{\inserttitle}
	\framesubtitle{\insertshortsubtitle: Introduction}
	
	\begin{itemize}
		\item A text is structured in paragraphs, sentences, words and characters.
		\item Some tasks use this structure.
		\item Many tasks rely on words as their processing unit.
		\item Preprocessing phase:
		\begin{itemize}
			\item Segmenting text;
			\item Filtering out unnecessary data;
			\item Normalizing words to limit their variations.
		\end{itemize}
	\end{itemize}

\end{frame}


\begin{frame}
	\frametitle{\inserttitle}
	\framesubtitle{\insertshortsubtitle: Plan}
	
	\begin{multicols}{2}
	%	\small
	\tableofcontents
	\end{multicols}
\end{frame}

%===================================================================================
\section{Characters}
%===================================================================================

\begin{frame}
	\frametitle{\insertshortsubtitle}
	\framesubtitle{\insertsection}

	\begin{itemize}
		\item \optword{Regular expressions}
		\begin{itemize}
			\item used to search for strings in text.
			\item recognizes languages of type 3 (regular languages) in Chomsky hierarchy.
			\item useful for lexical analysis (word tokenization).
		\end{itemize}
		\item \optword{Edit distance}
		\begin{itemize}
			\item used to measure the difference between two strings.
			\item useful for fuzzy search.
		\end{itemize}
	\end{itemize}

\end{frame}

\subsection{Regular Expressions}

\begin{frame}
	\frametitle{\insertshortsubtitle: \insertsection}
	\framesubtitle{\insertsubsection: Syntax (1)}

	\SetTblrInner{rowsep=0pt,colsep=1pt}
	\begin{tblr}{
			colspec = {p{.1\textwidth}p{.41\textwidth}p{.45\textwidth}},
			row{odd} = {lightblue, font=\small},
			row{even} = {lightyellow, font=\small},
			row{1} = {darkblue, font=\bfseries},
		} 
		\textcolor{white}{RE} & \textcolor{white}{Meaning} & \textcolor{white}{Example} \\
		
		. & any character & \keyword{beg.n}: I \expword{begun} at the \expword{begin}ning. \\
		
		 \empty [aeuio] & specific characters & \keyword{[Cc][au]t}: \expword{Cat}s are \expword{cut}e. \\
		 
		\empty [a-e] & character range & \keyword{[A-Z]..}: \expword{I s}aw \expword{Kar}im. \\
		
		\empty [\textasciicircum aeuio] & exclude characters & \keyword{[\textasciicircum A-Z]a.}: I \expword{saw} Karim. \\
		
		c? & one or zero & \keyword{colou?r}: It is \expword{colour} or \expword{color}. \\
		
		c* & zero or more & \keyword{Ye*s} : \expword{Ys}! \expword{Yes}! \expword{Yeeees}! \\
		
		c+ & one or more & \keyword{Ye+s} : Ys! \expword{Yes}! \expword{Yeeees}! \\
		
		c\{n\} & n occurrences & \keyword{Ye\{3\}s} : Ys! Yes! Yees! \expword{Yeees}! \\
		
		c\{n,m\} & between n and m occurrences & \keyword{Ye\{1,2\}s} : Ys! \expword{Yes}! \expword{Yees}! Yeees! \\
		
		c\{n,\} & at least n occurrences & \keyword{Ye\{2,\}s} : Ys! Yes! \expword{Yees}! \expword{Yeees}! \\
		
		c\{,m\} & at most m occurrences & \keyword{Ye\{,2\}s} : \expword{Ys}! \expword{Yes}! \expword{Yees}! Yeees! \\
		
	\end{tblr}

\end{frame}

\begin{frame}
	\frametitle{\insertshortsubtitle: \insertsection}
	\framesubtitle{\insertsubsection: Syntax (2)}

	\begin{minipage}{.6\textwidth}
	\begin{itemize}
		\item Grouping: \keyword{( )}
		\begin{itemize}
			\item E.g., \expword{/([Pp]retty, )+/}: \expword{Pretty, pretty, }pretty good.
		\end{itemize}
		\item Disjunction: \keyword{\textbar}
		\begin{itemize}
			\item E.g., \expword{/continu(e\textbar ation\textbar al(ly)?)/}
		\end{itemize}
		\item Text beginning: \keyword{\textasciicircum}
		\begin{itemize}
			\item E.g., \expword{/\textasciicircum K/}:  \expword{K}ill Karim.
		\end{itemize}
		\item Text ending: \keyword{\$}
		\begin{itemize}
			\item E.g., \expword{/\textbackslash .[\textasciicircum .]+\$/} :  file.tar\expword{.gz}
		\end{itemize}
		\item To capture a group: \keyword{\$n} or \keyword{\textbackslash n} where \expword{n} is the group's position.
		\begin{itemize}
			\item E.g., \expword{/(.*)(ual\textbar uation)\$/\textbackslash 1ue/}.
		\end{itemize}
	\end{itemize}
	\end{minipage}
	\begin{minipage}{.38\textwidth}
	\begin{tblr}{
			colspec = {p{.2\textwidth}p{.6\textwidth}},
			row{odd} = {lightblue},
			row{even} = {lightyellow},
			row{1} = {darkblue},
		} 
	
		\textcolor{white}{RE} & \textcolor{white}{Equivalence} \\
		
		\textbackslash d & [0-9] \\
		\textbackslash D & [\textasciicircum 0-9] \\
		\textbackslash w & [a-zA-Z0-9\_] \\
		\textbackslash W & [\textasciicircum \textbackslash w] \\
		\textbackslash s & [ \textbackslash r\textbackslash t\textbackslash n\textbackslash f] \\
		\textbackslash S & [\textasciicircum \textbackslash s] \\
	\end{tblr}
	
	\end{minipage}

\end{frame}

\begin{frame}
	\frametitle{\insertshortsubtitle: \insertsection}
	\framesubtitle{\insertsubsection: Use case}

	\begin{itemize}
		\item Text editors use REs for search and replacement.
		\item Most programming languages introduce a mechanism to use REs.
		\item Data Extraction: E.g., \expword{extract emails and phone numbers from blogs and social media}.
		\item \url{https://github.com/kariminf/aruudy}
		\begin{itemize}
			\item \textit{One of my crazy projects}.
			\item Arabic poems meter detection.
			\item Pros: readable rules.
			\item Cons: many passes to process one verse.
		\end{itemize}
	\end{itemize}

\end{frame}


\begin{frame}
	\frametitle{\insertshortsubtitle: \insertsection}
	\framesubtitle{\insertsubsection: Exercise}
	
	\begin{itemize}
		\item Find the regular expression which searches all conjugations of the verb ``\textbf{go}" (whithout auxilary verb). 
		\item Given a log, we want to list all lines starting with ``\textbf{Error}"; containing a number (which starts with a digit other than zero, has 3 to 5 consecutive zeros and ends with a digit other than zero); and ending with ``\textbf{...}"
		\item Find a regular expression which searches for words containing letters ``l", ``i" and ``n" in this order. The beginning of the word can be in uppercase, the rest in lowercase. For example, \expword{lion, Linux, violin, absolution, Aladdin, ...}
	\end{itemize}
	
\end{frame}

\begin{frame}
	\frametitle{\insertshortsubtitle: \insertsection}
	\framesubtitle{\insertsubsection: Some humor}

	\begin{center}
		\vgraphpage{humor/humor-regex.jpg}
	\end{center}

\end{frame}

\subsection{Edit distance}

\begin{frame}
	\frametitle{\insertshortsubtitle: \insertsection}
	\framesubtitle{\insertsubsection: Edit operations}

	\begin{itemize}
		\item \optword{Insertion}: inserting a character to a string.\\
		$uv \rightarrow uxv \,/\, u, v \in X^*;\, uv \in X^+;\, x \in X$
		\begin{itemize}
			\item \expword{case $ \rightarrow $ casse, entraînement $ \rightarrow $ entraînnement }
		\end{itemize}
		
		\item \optword{Deletion}: deleting a character from a string.\\
		$uxv \rightarrow uv \,/\, u, v \in X^*;\, uv \in X^+;\, x \in X$
		\begin{itemize}
			\item \expword{future $ \rightarrow $ futur, héros $ \rightarrow $ héro}
		\end{itemize}
		
		\item \optword{Substitution}: substituting a character by another.\\
		$uxv \rightarrow uyv \,/\, u, v \in X^*;\, x, y \in X;\, x \ne y$
		\begin{itemize}
			\item \expword{design $ \rightarrow $ dezign, croient $ \rightarrow $ croyent }
		\end{itemize}
		
		\item \optword{Transposition}: changing two characters position.\\
		$uxwyv \rightarrow uywxv \,/\, u, v, w \in X^*;\, x, y \in X;\, x \ne y$
		\begin{itemize}
			\item \expword{play $ \rightarrow $ paly, cueillir $ \rightarrow $ ceuillir}
		\end{itemize}
	\end{itemize}

\end{frame}

\begin{frame}
	\frametitle{\insertshortsubtitle: \insertsection}
	\framesubtitle{\insertsubsection: Algorithms}

	\begin{itemize}
		\item \optword{Hamming  distance}: allows only substitution. Strings must have the same size. 
		\item \optword{The longest common subsequence (LCS)}: allows only deletion and insertion.
		\item \optword{Levenshtein distance}: allows insertion, deletion and substitution.
		\item \optword{Jaro distance}: allows only transposition.
		\item \optword{Damerau–Levenshtein distance}: allows insertion, deletion, substitution end transposition between two  adjacent characters.
	\end{itemize}

\end{frame}


\begin{frame}
	\frametitle{\insertshortsubtitle: \insertsection}
	\framesubtitle{\insertsubsection: Hamming distance}
	
	\begin{itemize}
		\item X: source string of length $n$.
		\item Y: destination string of length $n$.
	\end{itemize}
	
	\[
	D_{i} = D_{i-1} + 
	\begin{cases}
	0 \text{ if } X_i = Y_i\\
	1 \text{ otherwise}
	\end{cases}
	\]
	
	\begin{exampleblock}{Examples of Hamming distance}
		\begin{itemize}
			\item D(100100, 101001) = 3
			\item D(abdelkrim, abderahim) = 3
		\end{itemize}
	\end{exampleblock}
	
\end{frame}

\begin{frame}
	\frametitle{\insertshortsubtitle: \insertsection}
	\framesubtitle{\insertsubsection: Levenshtein distance}

	\begin{itemize}
		\item X: source string of length $n$.
		\item Y: destination string of length $m$.
		\item D: a matrix; where D[i, j] is the edit distance between the substrings X[1..i] and Y[1..j].
		\item To calculate $D[n, m]$, dynamic programming is used.
		\item $D[0, 0] = 0$.
	\end{itemize}
	
	\[
	D[i, j] = \min 
	\begin{cases}
	D[i - 1, j] + 1 \text{ //Deletion}\\
	D[i, j-1] + 1 \text{ //Insertion}\\
	D[i-1, j-1] + \begin{cases}
	2 & \text{ if } X[i] \ne Y[j] \\
	0 & \text{ otherwise}
	\end{cases}
	\end{cases}
	\]

\end{frame}

\begin{frame}
	\frametitle{\insertshortsubtitle: \insertsection}
	\framesubtitle{\insertsubsection: Levenshtein distance (Example)}

	\begin{exampleblock}{Example of Levenshtein distance calculation}
		\centering\scriptsize
		\begin{tblr}{
				colspec = {|Q[c,m,wd=3em]|Q[c,m,wd=3em]|c|c|c|c|c|c|c|c|},
				row{1} = {bg=darkblue, fg=white, font=\bfseries},
				column{1} = {bg=darkblue, fg=white, font=\bfseries},
				stretch=1.5,
				colsep=2pt,
			} 	
			
			\hline
			& \# & i &\bfseries m &\bfseries m & a & t & u & r & e \\
			\hline
			\bfseries \# & \SetCell{bg=green!70} 0 & $ \leftarrow $ 1 & $ \leftarrow $ 2 & $ \leftarrow $ 3 & $ \leftarrow $ 4 & $ \leftarrow $ 5 & $ \leftarrow $ 6 & $ \leftarrow $ 7 & $ \leftarrow $ 8\\
			\hline
			\bfseries a & \SetCell{bg=green!70} $ \uparrow $ 1 & $ \nwarrow\leftarrow\uparrow $ 2 & $ \nwarrow\leftarrow\uparrow $ 3 & $ \nwarrow\leftarrow\uparrow $ 4 & $ \nwarrow $ 3 & $ \leftarrow $ 4 & $ \leftarrow $ 5 & $ \leftarrow $ 6 & $ \leftarrow $ 7 \\
			\hline
			\bfseries m & \SetCell{bg=green!70} $ \uparrow $ 2 & $ \nwarrow\leftarrow\uparrow $ 3 & $\nwarrow $ 2 & $\nwarrow\leftarrow $ 3 & $\leftarrow\uparrow $ 4 & $\nwarrow\leftarrow\uparrow $ 5 & $\nwarrow\leftarrow\uparrow $ 6 & $\nwarrow\leftarrow\uparrow $ 7 & $\nwarrow\leftarrow\uparrow $ 8\\
			\hline
			\bfseries i & $ \uparrow $ 3 & \SetCell{bg=green!70} $ \nwarrow $ 2 & \SetCell{bg=green!70} $\leftarrow\uparrow $ 3 & \SetCell{bg=green!70} $\nwarrow\leftarrow\uparrow $ 4 & \SetCell{bg=green!70} $\nwarrow\leftarrow\uparrow $ 5 & \SetCell{bg=green!70} $\nwarrow\leftarrow\uparrow $ 6 & \SetCell{bg=green!70} $\nwarrow\leftarrow\uparrow $ 7 & $\nwarrow\leftarrow\uparrow $ 8 & $\nwarrow\leftarrow\uparrow $ 9\\
			\hline
			\bfseries n & $ \uparrow $ 4 & $ \uparrow $ 3 & $\nwarrow\leftarrow\uparrow $ 4 & $\nwarrow\leftarrow\uparrow $ 5 & $\nwarrow\leftarrow\uparrow $ 6 & $\nwarrow\leftarrow\uparrow $ 7 & $\nwarrow\leftarrow\uparrow $ 8 & \SetCell{bg=green!70} $\nwarrow\leftarrow\uparrow $ 9 & $\nwarrow\leftarrow\uparrow $ 10\\
			\hline
			\bfseries e & $ \uparrow $ 5 & $ \uparrow $ 4 & $\nwarrow\leftarrow\uparrow $ 5 & $\nwarrow\leftarrow\uparrow $ 6 & $\nwarrow\leftarrow\uparrow $ 7 & $\nwarrow\leftarrow\uparrow $ 8 & $\nwarrow\leftarrow\uparrow $ 9 & $\nwarrow\leftarrow\uparrow $ 10 & \SetCell{bg=green!70} $\nwarrow $ 9\\
			\hline
		\end{tblr}
	\end{exampleblock}

\end{frame}

\begin{frame}
	\frametitle{\insertshortsubtitle: \insertsection}
	\framesubtitle{\insertsubsection: Some applications}

	\begin{itemize}
		\item \optword {file revision}: E.g., the UNIX \expword{diff} command which compares two files.
		\item \optword {spell checking}: suggesting possible corrections for a spell error (E.g. \expword{Hunspell}).
		\item \optword {plagiarism detection}: Here, words are used instead of characters.
		\item \optword {spam filtering}: Sometimes spammers intentionally use spell errors to deceive spam detection tools.
		\item \optword {bioinformatics}: similarity between two DNA sequences.
	\end{itemize}

\end{frame}

\begin{frame}
	\frametitle{\insertshortsubtitle: \insertsection}
	\framesubtitle{\insertsubsection: Exercise}
	
	\begin{itemize}
		\item In order to write the word "magnificent", the wrong word: "magnecifint" was written instead. Indicate each edit operation position(s) according to the correct word (If the operation does not exist, write 0. Transposition and substitution are prioritized, i.e. they must not e considered as a combination of insertion/deletion operations):
		\begin{tabular}{|lll|}
			\hline 
			Insertion : ............ & & Substitution  : ............ \\
			Deletion  : ............ & & Transposition : ............ \\
			\hline
		\end{tabular}
		
		\item Calculate Hamming and Levenstein distances of the two words: ``tray" et ``tary".
		Indicate different edit operation for each distance. Redo the same thing with Lavenstein distance (substitution weight = 1).
	\end{itemize}
	
\end{frame}

\begin{frame}
	\frametitle{\insertshortsubtitle: \insertsection}
	\framesubtitle{\insertsubsection: Some humor}

	\begin{center}
		\vgraphpage[.4\textheight]{humor/humor-spell.jpg}
		\vgraphpage[.4\textheight]{humor/humor-spell1.jpg}
		\vgraphpage[.4\textheight]{humor/humor-spell2.jpg}
	\end{center}

\end{frame}

%===================================================================================
\section{Text Segmentation}
%===================================================================================

\begin{frame}
	\frametitle{\insertshortsubtitle}
	\framesubtitle{\insertsection}

	\begin{itemize}
		\item New paragraphs are marked by a line break, or using HTML tag \keyword{\textless p\textgreater}.
		\begin{itemize}
			\item Using tools to extract text from PDFs, many line breaks may occur. 
		\end{itemize}
		\item New sentences are marked by a period (or another marker).
		\begin{itemize}
			\item Periods are not always used to mark sentence ending.
		\end{itemize}
		\item Words are separates by blanks.
		\begin{itemize}
			\item Not always!
		\end{itemize}
	\end{itemize}

\end{frame}

\subsection{Sentence boundary disambiguation}

\begin{frame}
	\frametitle{\insertshortsubtitle: \insertsection}
	\framesubtitle{\insertsubsection: Problematic} 

	\begin{itemize}
		\item \expword{/[.?!]/} is a simple RE to delimit sentences (French, English, etc.).
		\item Periods are used in numbers: \expword{123,456.78 (American style) 123.456,78 (European style)}
		\item Abbreviations contain periods.
		\item Long sentences are hard to process; it will be better to split them.
		\item Quotations.
		\item Some languages, such as Thai, do not use markers to delimit sentences.
	\end{itemize}

\end{frame}

\begin{frame}
	\frametitle{\insertshortsubtitle: \insertsection}
	\framesubtitle{\insertsubsection: Features} 

	\begin{itemize}
		\item \optword{Letter case}: sentences start with uppercase.
		\item \optword{Parts of speech (PoS)}: PoS can be used to decide a sentence boundary.
		\item \optword{Word length}: abbreviations are usually short.
		\item \optword{Prefixes and suffixes}: words having infixes are less probable to be abbreviations.
		\item \optword{Abbreviations classes} 
		\item \optword{Proper nouns}: proper nouns always start with an uppercase; they can occur in the middle of a sentence.
	\end{itemize}

\end{frame}

\begin{frame}
	\frametitle{\insertshortsubtitle: \insertsection}
	\framesubtitle{\insertsubsection: Some humor} 


	\begin{center}
		\vgraphpage{humor/humor-segmentation.jpg}
	\end{center}

\end{frame}

\subsection{Word Tokenization}

\begin{frame}
	\frametitle{\insertshortsubtitle: \insertsection}
	\framesubtitle{\insertsubsection: Problematic} 

	\begin{itemize}
		\item \expword{/[ ]+/} is a RE used to separate words (Arabic, French, English, etc.)
		\item Some languages, such as Japanese, do not use markers to separate words (\expword{\begin{CJK}{UTF8}{min}今年は本当に忙しかったです。\end{CJK}}).
		\item Compound words: open (``\expword{wake up}"), closed (``\expword{makeup}"; German: ``\expword{Lebensversicherung}: life insurance"; Arabic: ``\expword{\RL{y_htbrwnhm}}: they test them") or hyphenated (``\expword{well-prepared}"; French: ``\expword{va-t-il}: will he"; ``\expword{c-à-dire}: which means").
		\item Confusion between characters: apostrophe is used for quotation or contraction (\expword{She's, J'ai}).
		\item Multiple words expressions: \expword{numerical expressions such as dates}.
	\end{itemize}

\end{frame}

\begin{frame}
	\frametitle{\insertshortsubtitle: \insertsection}
	\framesubtitle{\insertsubsection: Approaches}

	\begin{itemize}
		\item \optword{Rules}: using regular expressions (REs).
		\begin{itemize}
			\item \url{https://www.nltk.org/api/nltk.tokenize.html}
			\item \url{https://nlp.stanford.edu/software/tokenizer.shtml}
			\item \url{https://spacy.io/}
			\item \url{https://github.com/kariminf/jslingua}
			\item \url{https://github.com/linuxscout/pyarabic}
		\end{itemize}
		\item \optword{Statistics}: using language models to calculate the probability of a character being the marker of word boundary.
		\begin{itemize}
			\item \url{https://nlp.stanford.edu/software/segmenter.html}
			\item \url{https://opennlp.apache.org/}
		\end{itemize}
	\end{itemize}

\end{frame}

\begin{frame}
	\frametitle{\insertshortsubtitle: \insertsection}
	\framesubtitle{\insertsubsection: Some humor}

	\begin{center}
		\vgraphpage{humor/humor-tokenization.png}
	\end{center}

\end{frame}


%===================================================================================
\section{Text normalization and filtering}
%===================================================================================

\begin{frame}
	\frametitle{\insertshortsubtitle}
	\framesubtitle{\insertsection}

	\begin{itemize}
		\item \optword{Text normalization}
		\begin{itemize}
			\item \textbf{PROBLEMATIC}: Texts can contain different variations of the same term. 
			Also, some tasks, such as information retrieval, do not need the exact text content.
			\item \textbf{SOLUTION}: transform a text to its canonical form.
			\item \url{https://www.kaggle.com/c/text-normalization-challenge-english-language}
		\end{itemize}
		\item \optword{Text filtering}
		\begin{itemize}
			\item \textbf{PROBLEMATIC}: Texts can contain some non useful characters, words or expressions.
			\item \textbf{SOLUTION}: deletion.
		\end{itemize}
	\end{itemize}

\end{frame}

\subsection{Text normalization}

\begin{frame}
	\frametitle{\insertshortsubtitle: \insertsection}
	\framesubtitle{\insertsubsection (1)}

	\begin{itemize}
		\item Acronyms and abbreviations;
		\begin{itemize}
			\item standard form: \expword{US \textrightarrow\ USA, U.S.A. \textrightarrow\ USA}
			\item long version: \expword{Mr. \textrightarrow\ Mister}
		\end{itemize}
		
		\item Using the same formatting for values such as dates and numbers;
		\begin{itemize}
			\item conversion to textual form: \expword{1205 DZD \textrightarrow\ One thousand, two hundreds, five Algerian dinars}.
			\item specific format: \expword{12 janvier 1986, 12.01.86, January 12th, 1986 \textrightarrow\ 1986-01-12}.
			\item replacement by type: \expword{January 12th, 1986 \textrightarrow\ DATE, kariminfo0@gmail.com \textrightarrow\ EMAIL}.
		\end{itemize}
		
		\item Transform uppercase to lowercase;
		\begin{itemize}
			\item \expword{Text \textrightarrow\ text}
			\item Sometimes, letters case must not transformed, such as in proper nouns (\expword{Will})
		\end{itemize}
	\end{itemize}

\end{frame}

\begin{frame}
	\frametitle{\insertshortsubtitle: \insertsection}
	\framesubtitle{\insertsubsection (2)}

	\begin{itemize}
		
		\item Contractions;
		\begin{itemize}
			\item \expword{y'll \textrightarrow\ you all, s'il \textrightarrow\ si il}
		\end{itemize}
		
		\item Diacritics;
		\begin{itemize}
			\item unaccentuation:  \expword{système \textrightarrow\ systeme}
			\item unvocalization:  \expword{\RL{yadrusu} \textrightarrow\ \RL{ydrs}}. Unless diacritics must be used (poems)
		\end{itemize}
	
		\item Encoding.
		\begin{itemize}
			\item The same encoding must be used for processing
		\end{itemize}
	
	\end{itemize}

\end{frame}

\subsection{Text filtering}

\begin{frame}
	\frametitle{\insertshortsubtitle: \insertsection}
	\framesubtitle{\insertsubsection}

	\begin{itemize}
		\item Special characters such as non printable characters;
		\begin{itemize}
			\item they can lead to erroneous processing. 
		\end{itemize}
		\item Keywords in textual formats;
		\begin{itemize}
			\item Tags: HTML, XML, etc. 
		\end{itemize}
		\item \keyword{stop words}: insignificant words such as prepositions, articles and pronouns.
		\begin{itemize}
			\item In information retrieval, there is no need to index them.
			\item In automatic text summarization, these words can alter sentence scores.
		\end{itemize}
	\end{itemize}

\end{frame}

\begin{frame}
	\frametitle{\insertshortsubtitle: \insertsection}
	\framesubtitle{\insertsubsection: Some humor}

	\begin{center}
		\vgraphpage{humor/humor-stopwords.jpg}
	\end{center}

\end{frame}

%===================================================================================
\section{Morphology}
%===================================================================================

\begin{frame}
	\frametitle{\insertshortsubtitle}
	\framesubtitle{\insertsection}

	\begin{itemize}
		\item Some languages admit words forming using flexion (e.g. \expword{conjugation}) and derivation (e.g. \expword{nominalization}).
		\item The most used method is affixation.
		\item Automatizing this task can help other tasks, such as natural language generation (NLG).
		\item Also, the inverse task can help other tasks, such as natural language comprehension (NLU).
	\end{itemize}

\end{frame}

\subsection{Word formation}

\begin{frame}
	\frametitle{\insertshortsubtitle: \insertsection}
	\framesubtitle{\insertsubsection}

	\begin{itemize}
		\item \optword{Flexion} : morphological variation of a word according to grammatical features (number, gender, etc.).
		\begin{itemize}
			\item Verbs conjugation.
			\item Nouns, pronouns, adjectives, adverbs, and articles declension. 
			E.g. \expword{\RL{qi.t.t} \textrightarrow\ \RL{qi.ta.t}}.
		\end{itemize}
		\item \optword{Derivation}: morphological variation of a word to create a new lexeme or to change category.
		\begin{itemize}
			\item Create a new lexeme: \expword{couper \textrightarrow\ découper, \RL{`ml} \textrightarrow\ \RL{ist`ml}}.
			\item Change category: nominalization (\expword{classer \textrightarrow\ classement, classeur ; \RL{darasa}. \textrightarrow\ \RL{darsuN, madrasaTuN, mudarrisuN, dArisuN}}), adjective (\expword{fatiguer \textrightarrow\ fatigant}).
		\end{itemize}
	\end{itemize}

\end{frame}

\begin{frame}
	\frametitle{\insertshortsubtitle: \insertsection}
	\framesubtitle{\insertsubsection: Example (Automatic conjugation)}

	\begin{itemize}
		\item \optword{Databases}:
		\begin{itemize}
			\item storing verbs conjugations in a database.
		\end{itemize}
		\item \optword{Template}:
		\begin{itemize}
			\item storing conjugations of certain verbs as models and use a list to index similar verbs.
			\item E.g., \expword{French verbs conjugation}.
		\end{itemize}
		\item \optword{Rules}:
		\begin{itemize}
			\item using IF-ELSE rules and regular expression.
			\item E.g., Arabic, English, French and Japanese verbs conjugation: \url{https://github.com/kariminf/jslingua}
			\item E.g., Arabic verbs conjugation: \url{https://github.com/linuxscout/qutrub}
		\end{itemize}
	\end{itemize}

\end{frame}

\begin{frame}
	\frametitle{\insertshortsubtitle: \insertsection}
	\framesubtitle{\insertsubsection: Some humor}

	\begin{center}
		\vgraphpage{humor/humor-formation.jpg}
	\end{center}

\end{frame}

\subsection{Form reducing}

\begin{frame}
	\frametitle{\insertshortsubtitle: \insertsection}
	\framesubtitle{\insertsubsection: Stemming}

	\begin{itemize}
		\item Affixes deletion.
		\item Result: \keyword{stem}.
		\item E.g., \expword{chercher \textrightarrow\ cherch}, \expword{arguing \textrightarrow\ argu}.
		\item \optword{Database}: storing all terms and their stems in a table.
		\item \optword{Statistics}: using a language model (N-Gram) to estimate truncation position.
		\item \optword{Rules}:
		\begin{itemize}
			\item Porter algorithm \cite{1980-porter}.
			\item \url{https://github.com/assem-ch/arabicstemmer}
			\item \url{https://snowballstem.org/algorithms/}
		\end{itemize}
	\end{itemize}

\end{frame}

\begin{frame}
	\frametitle{\insertshortsubtitle: \insertsection}
	\framesubtitle{\insertsubsection: Stemming (E.g. Porter algorithm)}

	\begin{itemize}
		\item A set of rules: condition/action.
		\item \url{https://snowballstem.org/}: a framework to create stemmers
		\item \optword{Condition on the stem}:
		\begin{itemize}
			\item length, ending, if it contains vowels, etc.
			\item E.g., \expword{(*v*) Y \textrightarrow\ I: happy \textrightarrow\ happi, sky \textrightarrow\ sky}
		\end{itemize}
		\item \optword{Condition on the affix}:
		\begin{itemize}
			\item For Porter stemmer, there is only suffixes 
			\item E.g., \expword{SSES \textrightarrow\ SS, ATIONAL \textrightarrow\ ATE}
		\end{itemize}
		\item \optword{Condition on the rules}:
		\begin{itemize}
			\item if a rule was executed, others must not.
		\end{itemize}
	\end{itemize}

\end{frame}

\begin{frame}
	\frametitle{\insertshortsubtitle: \insertsection}
	\framesubtitle{\insertsubsection: Lemmatization}

	\begin{itemize}
		\item Search canonical form of a word.
		\item Result: \keyword{lemma}.
		\item E.g., \expword{comprennent \textrightarrow comprendre, better \textrightarrow good}.
		\item Word's context is used. E.g. \expword{saw \textrightarrow\ (V) see ou (N) saw}.
		\item \optword{Lexical bases}:
		\begin{itemize}
			\item \url{https://www.nltk.org/api/nltk.stem.html} (NLTK WordNetLemmatizer)
			\item \url{https://github.com/sloria/textblob}
			\item \url{https://spacy.io/}
		\end{itemize}
		\item \optword{Machine learning}:
		\begin{itemize}
			\item \url{https://opennlp.apache.org/}
		\end{itemize}
	\end{itemize}

\end{frame}

\begin{frame}[fragile]
	\frametitle{\insertshortsubtitle: \insertsection}
	\framesubtitle{\insertsubsection: Lemmatization (E.g. Morphy)}

	\begin{block}{Wordnet's "morphy" Lemmatization}
		\footnotesize
		\begin{algorithm}[H]
			\vskip-3pt
			\KwData{word, PoS}
			\KwResult{list of possible lemmas}
			
			\If{word $ \in $  exceptions\_list[PoS]}{
				\Return search\_in\_dictionary(\{word\} $ \cup $ exceptions\_list[PoS])\;
			}
			
			forms = \{word\}\;
			
			\While{forms $ \ne \emptyset $}{
				forms = delete\_affixes(forms, PoS)\;
				
				results = search\_in\_dictionary(\{word\} $ \cup $ forms)\;
				
				\lIf{results $ \ne \emptyset $ }{
					\Return results
				}
			}
			
			\Return $ \emptyset $\;
	
		\end{algorithm}
	\end{block}

\end{frame}

\begin{frame}
	\frametitle{\insertshortsubtitle: \insertsection}
	\framesubtitle{\insertsubsection: Some humor}

	\begin{center}
		\vgraphpage{humor/humor-stemming.jpeg}
	\end{center}

\end{frame}

\insertbibliography{NLP03}{*}

\end{document}

